\section{The Entropy Condition}
\label{sec:entropy-condition}

Corollary~\ref{cor:nonuniqueness} showed that the weak formulation alone
does not determine a unique solution. This section introduces an
additional condition, motivated by the vanishing-viscosity limit made
precise in Section~\ref{sec:extensions}, that selects the physically
correct weak solution among any competing candidates.

\subsection{The Lax entropy condition}
\label{subsec:lax-condition}

\begin{definition}[Lax entropy condition]
\label{def:lax-entropy}
Let $u$ be a weak solution of a scalar conservation law $u_t+f(u)_x=0$
with $f$ convex, having a jump discontinuity along a $C^1$ curve $x=s(t)$
with left and right limits $u_\ell(t)$, $u_r(t)$. The jump satisfies the
\emph{Lax entropy condition} if
\begin{equation}
  \label{eq:lax-condition}
  f'(u_\ell(t)) > s'(t) > f'(u_r(t))
\end{equation}
for all $t$. A weak solution is an \emph{entropy solution} if every jump
discontinuity it contains satisfies \eqref{eq:lax-condition}.
\end{definition}

For Burgers' equation, $f'(u)=u$, so \eqref{eq:lax-condition} together
with the Rankine--Hugoniot speed $s'=(u_\ell+u_r)/2$ reduces to
\begin{equation}
  \label{eq:lax-burgers}
  u_\ell > \frac{u_\ell+u_r}{2} > u_r
  \quad \Longleftrightarrow \quad
  u_\ell > u_r.
\end{equation}

\begin{remark}
Condition \eqref{eq:lax-condition} says the characteristics on either side
of the shock must run \emph{into} it as $t$ increases, rather than out of
it: information is being destroyed at the shock (consistent with physical
irreversibility), not created. This is the precise sense in which a shock
differs from a rarefaction, where characteristics instead fan \emph{out}
of a single point (Section~\ref{subsec:nonuniq-rarefaction}).
\end{remark}

\subsection{Application to the Riemann problem of Section~\ref{sec:nonuniqueness}}
\label{subsec:entropy-application}

\begin{proposition}
\label{prop:entropy-selects-B}
For the Riemann data \eqref{eq:riemann-data} ($u_\ell=0<u_r=1$), the
shock solution $u_A$ of \eqref{eq:shock-solution} is not an entropy
solution, while the rarefaction solution $u_B$ of
\eqref{eq:rarefaction-solution} is.
\end{proposition}

\begin{proof}
For $u_A$, condition \eqref{eq:lax-burgers} requires $u_\ell>u_r$, but
here $u_\ell=0<1=u_r$, so \eqref{eq:lax-condition} fails: $u_A$ is not an
entropy solution. The function $u_B$ contains no jump discontinuity at
all (Proposition~\ref{prop:rarefaction-is-weak-solution} shows it is
continuous), so Definition~\ref{def:lax-entropy} imposes no condition on
it, and it is vacuously an entropy solution.
\end{proof}

Proposition~\ref{prop:entropy-selects-B} recovers, as a formal criterion,
the physical judgement made informally at the end of
Section~\ref{sec:nonuniqueness}: among the two weak solutions
constructed there, the entropy condition selects exactly the rarefaction
$u_B$ and rules out the backwards-running shock $u_A$.

\subsection{Uniqueness}
\label{subsec:entropy-uniqueness}

The entropy condition is not merely a rule of thumb for discarding
individual bad solutions: it in fact restores full well-posedness.

\begin{theorem}[Oleinik, uniqueness of entropy solutions; see
\cite{evans2010}, \S 3.4]
\label{thm:oleinik-uniqueness}
For $u_0 \in L^\infty(\mathbb{R})$ and $f$ uniformly convex, the initial
value problem $u_t+f(u)_x=0$, $u(\cdot,0)=u_0$, has at most one entropy
solution satisfying \eqref{eq:lax-condition} at every jump.
\end{theorem}

The proof of Theorem~\ref{thm:oleinik-uniqueness} (via an $L^1$-contraction
argument comparing two entropy solutions) is beyond the scope of this
paper and is not reproduced here; it is stated to record that the
selection principle of Definition~\ref{def:lax-entropy} is exactly strong
enough to recover uniqueness, and no stronger condition is needed.